% !TEX program = xelatex
\documentclass[10pt,a4paper]{article}

% ============================================================
% COMPILATION
% ============================================================
% Compile with XeLaTeX so that Telugu and Devanagari examples
% can be displayed correctly.

% ============================================================
% PACKAGES
% ============================================================
\usepackage[a4paper,margin=0.62in]{geometry}
\usepackage{graphicx}
\usepackage{booktabs}
\usepackage{multirow}
\usepackage{array}
\usepackage{amsmath}
\usepackage{float}
\usepackage{subcaption}
\usepackage{enumitem}
\usepackage{microtype}
\usepackage{hyperref}
\usepackage{xurl}
\usepackage{titlesec}
\usepackage{setspace}
\usepackage{fontspec}

% ============================================================
% FONTS
% ============================================================
\newfontfamily\telugufont{Noto Sans Telugu}
\newfontfamily\devanagarifont{Noto Sans Devanagari}

% ============================================================
% PAGE / FORMATTING
% ============================================================
\hypersetup{
    colorlinks=true,
    linkcolor=black,
    urlcolor=blue,
    citecolor=black
}

\setstretch{0.95}

\setlength{\parindent}{0pt}
\setlength{\parskip}{3pt}

\setlist[itemize]{
    nosep,
    leftmargin=*
}

\setlist[enumerate]{
    nosep,
    leftmargin=*
}

\titlespacing*{\section}{0pt}{6pt}{3pt}
\titlespacing*{\subsection}{0pt}{5pt}{2pt}

\titleformat{\section}
    {\large\bfseries}
    {\thesection}
    {0.5em}
    {}

\titleformat{\subsection}
    {\normalsize\bfseries}
    {\thesubsection}
    {0.5em}
    {}

% ============================================================
% FIGURE HELPER
% ============================================================
\newcommand{\safeimage}[2][\linewidth]{%
    \IfFileExists{#2}
    {\includegraphics[width=#1]{#2}}
    {\fbox{\parbox[c][1.7cm][c]{0.9\linewidth}
    {\centering Figure not uploaded}}}
}

% ============================================================
% TITLE
% ============================================================
\title{
    \vspace{-1.0cm}
    \textbf{Language Models and Agents: Individual Project Report}\\
}

\author{
    Korlapati Keerthana\\
    2023101121
}

\date{}

% ============================================================
% DOCUMENT
% ============================================================
\begin{document}

\maketitle
\vspace{-0.5cm}

% ============================================================
\section{Phase 1: Corpus Construction and Tokenization}
% ============================================================

\subsection{Language Selection}

The project requires one higher-resource non-English Indian language
and one lower-resource language from the specified list. Telugu was
selected as the higher-resource language and Nepali as the
lower-resource language.

Telugu was chosen as the higher resource language after finding substantially more text through publicly available sources. And also beacuse I can read Telugu, so can evaluate better. Nepali was selected from the permitted lower-resource language list and was processed independently from Telugu throughout the project.

\subsection{Corpus Construction}

For both languages, the corpus pipeline collected source documents,
removed invalid records and exact duplicates. Then all sources were merged, the data unicode normalized and deduplicated, and made into independent training, validation, and test sets. This entire corpus pipeline was done separately for both the languages.

\begin{table}[H]
\centering
\small
\caption{Corpus construction summary.}
\begin{tabular}{p{0.27\linewidth}cc}
\toprule
 & \textbf{Telugu} & \textbf{Nepali}\\
\midrule
Initial target & 600M tokens & 600M tokens\\
Final training corpus & 486.03M tokens & 500M tokens\\
Manual tokens available & 831.0M & 109.4M\\
Manual tokens selected & 150.0M & 109.4M\\
Manual fraction & 25.00\% & 18.24\%\\
Non-manual fraction & 75.00\% & 81.76\%\\
Validation tokens & 2.92M & 10M\\
Test tokens & 2.82M & 10M\\
\bottomrule
\end{tabular}
\end{table}

The Telugu corpus contained approximately 4.01M documents. The
collected material included Telugu news, Wikisource, Wikipedia,
Internet Archive, Gutenberg, thesis material, and IndicCorp.

The Nepali corpus contained approximately 1.01M documents, with
FineWeb2 and OnlineKhabar contributing majorly. Other sources
included Nepali Wikipedia, OCR material, parliamentary proceedings,
the Law Commission, and archival material.

The project requires at least 20\% of the final training tokens to
come from manual collection. Telugu had sufficient manual material,
allowing 150M manually collected tokens to be used. This gives a 25\%
manual fraction.

For Nepali, only approximately 109.4M manually collected tokens were
available. Consequently, 20\% of a 600M-token corpus could not be
reached using the available manual material. The selected corpus
therefore contains approximately 109.4M manual tokens, corresponding
to 18.24\%.

\subsection{Train/Test Leakage Prevention}

The validation and test data were selected separately from the
training data. For Telugu, training-record hashes were loaded before
sampling validation and test records. The final verification showed
zero overlap between training and validation, training and test, and
validation and test.

For the synthetic reasoning data, the generator also checks for
exact prompt overlap across splits. Train/validation templates and
test templates use separate template-ID families, while entity
overlap across splits is allowed. This is intentional because the
task is supposed to test the reasoning relation rather than whether
the model has seen a particular entity name before.

The generator also checks global prompt uniqueness and prevents test
template leakage into train or validation.

\subsection{Tokenizer Construction}

A separate BPE tokenizer was trained from scratch for each language.

\begin{table}[H]
\centering
\small
\caption{Final tokenizer statistics on held-out validation text.}
\begin{tabular}{lcc}
\toprule
 & \textbf{Telugu} & \textbf{Nepali}\\
\midrule
Vocabulary size & 6,000 & 8,000\\
Validation documents & 10,000 & 10,000\\
Validation tokens & 2,922,947 & 5,718,485\\
Fertility & 2.2417 & 1.6672\\
Unknown tokens & 0 & 0\\
Unknown-token rate & 0\% & 0\%\\
Tokens / character & 0.2830 & 0.2578\\
Characters / token & 3.53 & 3.88\\
Unique validation tokens & 4,387 & 4,597\\
Unused vocabulary entries & 1,613 & 3,403\\
\bottomrule
\end{tabular}
\end{table}

For Telugu, vocabulary sizes from 5K to 10K were compared. Increasing
the vocabulary reduced fertility, with fertility decreasing from
2.3850 at 5K to 1.9361 at 10K. The 6K vocabulary size was chosen as a
compute-aware compromise between token fragmentation and vocabulary
size.

\begin{table}[H]
\centering
\small
\caption{Telugu tokenizer vocabulary sweep.}
\begin{tabular}{lrrrrrr}
\toprule
Vocabulary & 5K & 6K & 7K & 8K & 9K & 10K\\
\midrule
Fertility
& 2.3850
& 2.2417
& 2.1358
& 2.0542
& 1.9911
& 1.9361\\
Tokens/char
& .3011
& .2830
& .2697
& .2594
& .2514
& .2444\\
\bottomrule
\end{tabular}
\end{table}

The Nepali tokenizer uses an 8K vocabulary and obtains zero unknown
tokens with fertility 1.6672. Large differences between sources were
observed. The OCR material has fertility of approximately 6.22,
compared to approximately 1.54 for OnlineKhabar and 1.67 for
FineWeb2. This can be interpreted as OCR noise producing more
fragmented token sequences.

\subsection{Token Frequency Statistics}

Token frequencies were also calculated on held-out validation text.
The Telugu plot below shows the most frequent tokens for the final
6K tokenizer, while the Nepali plot shows the corresponding
statistics for the final 8K tokenizer.

\begin{figure}[H]
\centering

\begin{subfigure}{0.48\linewidth}
    \centering
    \safeimage{figures/telugu_bpe_6000_top50_token_frequency.png}
    \caption{Telugu BPE, 6K vocabulary}
\end{subfigure}
\hfill
\begin{subfigure}{0.48\linewidth}
    \centering
    \safeimage{figures/nepali_validation_top_token_frequencies.png}
    \caption{Nepali BPE, 8K vocabulary}
\end{subfigure}

\caption{Top-token frequency plots on the validation sets.}
\end{figure}

The frequency distributions are concentrated around a relatively
small number of frequent tokens, while a large part of the vocabulary
is used less often. In the validation data, 4,387 of the 6,000 Telugu
tokens and 4,597 of the 8,000 Nepali tokens were observed.

\subsection{Tokenizer Examples}

The final tokenizer notebooks were also used to inspect how example
sentences are split into tokens and to check that decoding the token
IDs gives back the original sentence.

\begin{table}[H]
\centering
\scriptsize
\begin{tabular}{p{0.20\linewidth}p{0.72\linewidth}}
\toprule
\textbf{Language} & \textbf{Sentence $\rightarrow$ tokenized output}\\
\midrule
Telugu &
\telugufont
\textit{తెలుగు భారతదేశంలో మాట్లాడే ఒక ప్రధాన భాష. $\rightarrow$
'▁తెలుగు', '▁భారతదేశ', 'ంలో', '▁మాట్లా', 'డే', '▁ఒక', '▁ప్రధాన', '▁భాష', '.'
}
\\[8pt]
Nepali &
\devanagarifont
\textit{नेपाल प्राकृतिक सुन्दरताले भरिएको देश हो। $\rightarrow$ 'नेपाल', 'प्राकृतिक', 'सुन्दर', 'ताले', 'भरि', 'एको', 'देश', 'हो', '।'}
\\
\bottomrule
\end{tabular}
\caption{Tokenizer examples from the final tokenizer notebooks.}
\end{table}

The tokenizer checks also confirmed that the selected final
tokenizers produce zero unknown tokens on the held-out validation
sets.

% ============================================================
\section{Phase 2: Model Implementation, Pretraining, and Evaluation}
% ============================================================

\subsection{Model Configuration and Design Decisions}

For Phase 2, the following architecture was selected after examining
the tokenizer statistics and the approximately 25M-parameter
constraint given.

\begin{table}[H]
\centering
\small
\caption{Final Transformer configurations.}
\begin{tabular}{lcc}
\toprule
 & \textbf{Model H: Telugu} & \textbf{Model L: Nepali}\\
\midrule
Vocabulary size & 6,000 & 8,000\\
Model dimension & 448 & 448\\
Transformer layers & 10 & 11\\
Attention heads & 8 & 8\\
FFN dimension & 1,344 & 1,344\\
Context length & 512 & 512\\
Dropout & 0.10 & 0.10\\
Positional encoding & Learned absolute & Learned absolute\\
Normalization & Pre-norm & Pre-norm\\
Embedding tying & Yes & Yes\\
Trainable parameters & 23,042,432 & 25,950,848\\
\bottomrule
\end{tabular}
\end{table}

The same model dimension and number of attention heads were used for
both languages so that the comparison is more controlled. The number
of layers differs by one because the two models have different
vocabulary sizes, which also changes the number of parameters used by
the embedding and output layers.

Learned absolute positional embeddings were used because they give
the model direct information about token order and were also simple
to implement from scratch. Since the maximum sequence length is set
to 512, the position table is also limited to 512 positions. This
means the model has a fixed context length of 512 tokens.

Pre-norm Transformer blocks were used to make the training more
stable when stacking multiple Transformer blocks. The normalization
is applied before the attention and feed-forward sublayers.

The output embedding matrix is tied to the input embedding matrix.
This saves some parameters and helps keep the models within the
approximately 25M parameter limit.

\subsection{Transformer Implementation}

The models implement decoder-only causal self-attention directly
using PyTorch primitives. For an input tensor

\[
X \in \mathbb{R}^{B\times T\times d},
\]

queries, keys, and values are projected and reshaped from
$(B,T,d)$ to $(B,h,T,d_k)$. Attention is computed as

\[
\operatorname{Attention}(Q,K,V)
=
\operatorname{softmax}
\left(
\frac{QK^{T}}{\sqrt{d_k}}+M
\right)V,
\]

where $M$ is an upper-triangular causal mask.

The $1/\sqrt{d_k}$ factor keeps the query--key dot products from
becoming too large as the head dimension increases. Without this
scaling, the softmax can become too concentrated and make learning
harder. The causal mask makes sure that position $t$ cannot attend to
positions after $t$.

Each Transformer block contains causal multi-head attention, a GELU
feed-forward network, residual connections, layer normalization, and
dropout.

\subsection{Causal Mask Verification}

The causal mask was checked empirically in addition to being present
in the attention implementation.

For Telugu, the supplied verification changes only token position 6
in the input sequence. The model is run on the original and modified
sequences in evaluation mode. Since position 6 is in the future for
positions 0--5, the logits at positions 0--5 should not change. The
test checks that the maximum absolute logit difference at all these
positions is below $10^{-6}$.

The same-position and later-position logits are allowed to change,
since position 6 can see itself and position 7 can see position 6.
The Telugu test passed.

For Nepali, the supplied causal-mask test directly checks the
attention matrix and verifies that every future-token attention
weight is exactly zero. The test also passed.

Thus, the implementation was checked both through the attention
weights and, for Telugu, through the effect of changing a future
input token on earlier logits.

\subsection{Pretraining}

Both models were pretrained independently on their respective
corpora. AdamW, learning-rate warmup, cosine learning-rate
decay, validation evaluation, mixed precision, and checkpointing were
used.

\begin{table}[H]
\centering
\small
\caption{Pretraining configuration.}
\begin{tabular}{lcc}
\toprule
 & \textbf{Telugu} & \textbf{Nepali}\\
\midrule
Training tokens & 486.03M & 500M\\
Validation tokens & 2.92M & 10M\\
Batch size & 16 & 16\\
Peak learning rate & $3\times10^{-4}$ & $2\times10^{-4}$\\
Training steps & 59,214 & 61,000\\
Optimizer & AdamW & AdamW\\
Precision & FP16 & FP16\\
GPU & Tesla T4 & Tesla T4\\
\bottomrule
\end{tabular}
\end{table}

The final validation losses were 3.016 for Telugu and 3.834 for
Nepali. These correspond to validation perplexities of 20.40 and
46.25 respectively.

\begin{figure}[H]
\centering

\begin{subfigure}{0.47\linewidth}
    \centering
    \safeimage{figures/telugu_loss.png}
    \caption{Telugu}
\end{subfigure}
\hfill
\begin{subfigure}{0.47\linewidth}
    \centering
    \safeimage{figures/nepali_loss.png}
    \caption{Nepali}
\end{subfigure}

\caption{Training and validation loss curves for the two models.}
\end{figure}

\subsection{Language Modeling Evaluation}

\begin{table}[H]
\centering
\small
\caption{Held-out language-modeling results.}
\begin{tabular}{lcc}
\toprule
Metric & \textbf{Model H: Telugu} & \textbf{Model L: Nepali}\\
\midrule
Test cross-entropy & 3.124 & 3.409\\
Test perplexity & 22.73 & 30.23\\
Test BPB & 0.467 & 0.846\\
Evaluated target tokens & 2,824,704 & 9,999,872\\
Test UTF-8 bytes & 27.28M & 58.15M\\
\bottomrule
\end{tabular}
\end{table}

Telugu obtained lower cross-entropy, perplexity, and BPB.
The BPB difference is useful here because perplexity depends on the
tokenization scheme. The lower Telugu BPB suggests that the
difference is not only because the two models use different
vocabulary sizes.

The results also match some of the differences seen during Phase 1.
Both models were trained at approximately the same token scale, but
Telugu had more available manually collected material and used a
25\% manual selection. Nepali used the maximum available manual
material but reached only 18.24\%. Nepali also had an OCR component
with much higher tokenizer fertility.

% ============================================================
\subsection{Generation Quality and Diversity}
% ============================================================

Generation was evaluated from held-out prefixes using greedy decoding
and temperatures $0.5$, $1.0$, and $1.5$.

\begin{table}[H]
\centering
\scriptsize
\caption{Generation quality and diversity. BLEU is corpus BLEU-4;
chrF is chrF++; repetition is repeated 3-gram fraction.}
\begin{tabular}{llrrrrrr}
\toprule
Model & Setting & BLEU & chrF & ROUGE-L &
Rep. & D-1 & D-2\\
\midrule
\multirow{4}{*}{Telugu}
& Greedy & 8.14 & 22.28 & .137 & .543 & .302 & .405\\
& $T=0.5$ & 7.14 & 22.78 & .127 & .246 & .424 & .650\\
& $T=1.0$ & 5.41 & 22.81 & .097 & .009 & .727 & .969\\
& $T=1.5$ & .40 & 17.09 & .025 & .001 & .825 & .997\\
\midrule
\multirow{4}{*}{Nepali}
& Greedy & 3.43 & 16.66 & .135 & .712 & .116 & .226\\
& $T=0.5$ & 3.42 & 18.61 & .138 & .390 & .179 & .460\\
& $T=1.0$ & 1.62 & 18.92 & .102 & .014 & .323 & .894\\
& $T=1.5$ & .25 & 16.11 & .058 & .0002 & .412 & .981\\
\bottomrule
\end{tabular}
\end{table}

Higher temperatures create more diverse results, while lower
temperatures keep the output closer to the reference. There is a
direct trade-off here. Increasing temperature reduces repeated
n-grams and increases Distinct-1 and Distinct-2, but sufficiently
high temperatures reduce BLEU and ROUGE-L.

For Telugu, repeated 3-grams fall from 54.3\% under greedy decoding
to 0.9\% at $T=1.0$, while Distinct-2 increases from 0.405 to 0.969.
Nepali follows the same trend but starts from a more repetitive
greedy baseline.

\subsection{Actual Generated Samples}

The following are actual outputs from the held-out generation runs.
A single held-out prefix is shown for each language so that the
change with temperature can be seen more easily.

\begin{table}[H]
\centering
\scriptsize
\begin{tabular}{p{0.13\linewidth}p{0.79\linewidth}}
\toprule
\multicolumn{2}{c}{\textbf{Telugu}}\\
\midrule
Greedy &
\telugufont ఓడించి, ఆ పార్టీతో పొత్తు పెట్టుకుని, ఆ పార్టీతో పొత్తు
పెట్టుకుని, ఆ పార్టీతో పొత్తు పెట్టుకుని, ఆ పార్టీతో పొత్తు పెట్టుకుని,
ఆ పార్టీతో పొత్తు పెట్టుకుని, ఆ పార్టీతో పొత్తు
పెట్టుకుని, ఆ పార్టీతో పొత్తు పెట్టుకుని, ఆ పార్టీతో పొత్తు
\\
$T=0.5$ &
\telugufont గట్టిగా నిలదీసింది. ఆ పార్టీతో పాటు, ఆ పార్టీతో కూడా
విభేదాలు వచ్చాయి. ఆ పార్టీతో పాటు, ఆ పార్టీతోనూ విభేదాలు వచ్చాయి.
ఆ పార్టీతోనూ విభేదాలు వచ్చాయి. ఆ పార్టీతోనూ విభేదాలు వచ్చాయి.
\\
$T=1.0$ &
\telugufont ఆస్పత్రికి మార్చింది. అందువల్ల ఇక్కడ యాత్ర ముగించి
అక్కడకు వెళ్ళి అక్కడి నుంచి జారిపోవాలని నిర్ణయించుకుంది. కాని ఆ రోజు
తను పోలీసు బలగాలు చేతికి నిరసనగా వెళుతున్నప్పుడు కంట్రోల్ చేయడానికి
కాని ఎవరి మీద మంత్రులు ఉన్నారు.
\\
$T=1.5$ &
\telugufont జూన్ 16 టిడిపి ప్రతి రైడుతూ సలహాలు ఇవ్వాలనిపించింది.
ఆలోచన విధాన ప నారాయణ, వంశీ, మీహి, సాయికృష్ణ ఈ మార్పు ముఖ్యం...
సంజయమ్ము రవి, కళా ఆలోచనలు ఆయనకు ఉన్నాయి.
\\
\midrule
\multicolumn{2}{c}{\textbf{Nepali}}\\
\midrule
Greedy &
\devanagarifont मा रहेका मतदान केन्द्र हरुमा मतदान केन्द्र हरु
निर्माण गरिएको छ । अ त्त रिया मा रहेका मतदान केन्द्र हरुमा मतदान
केन्द्र हरु निर्माण गरिएको छ । अ त्त रिया मा रहेका मतदान केन्द्र हरुमा
मतदान केन्द्र हरु निर्माण गरिएको छ ।
\\
$T=0.5$ &
\devanagarifont मा रहेका मतदान केन्द्रमा मतदान अधिकृत र मतदान
अधिकृत लाई मतदान अधिकृत ले मतदान अधिकृत ले मतदान अधिकृत लाई मतदान
अधिकृत ले मतदान अधिकृत लाई मतदान अधिकृत ले मतदान अधिकृत ले मतदान
अधिकृत लाई मतदान अधिकृत ले
\\
$T=1.0$ &
\devanagarifont मा मात्र ३७ १ मतदाता रहेका छन् । अहिलेसम्म तिरेको
ओझेलमा परेको उक्त मतदान केन्द्रमा मणि पाल, टेक बहादुर विकको कक्षा १०
मा अध्ययनरत श्याम लाल चौधरी लगायतले मतदान गरेका छन् ।
\\
$T=1.5$ &
\devanagarifont यता नेपाली खेल पत्रकार अभियानमा स्तरीय वाद तयार
भएपनि तस् होस् : घनन्देही, ३१ साउन उजुरी रको अर्को फैसला वर्तमान
विषय घटना भएको ठा दद काउन्सिलका बसोवाससँग मिल्दो द्वासाईकेट
(बग एउटै भवन व्यवस्था पक्षले बनेको बाग बजार विकास
\\
\bottomrule
\end{tabular}
\caption{Actual generated continuations for one held-out prefix at
different temperatures.}
\end{table}

Qualitatively, the greedy outputs in both languages are quite
repetitive. At $T=0.5$, some variation is introduced but repeated
phrases are still common. At $T=1.0$, the outputs are more varied and
contain longer new continuations, although they are not always
fully coherent. At $T=1.5$, repetition is very low but the text
becomes less consistent and sometimes switches to unrelated topics or
phrasing. The outputs are generally in the correct script, but
language correctness and coherence are not consistent at higher
temperatures.

BLEU and ROUGE-L therefore need to be interpreted together with the
qualitative samples and diversity measures. The character-level chrF
metric is also useful because exact subword overlap is not required.

% ============================================================
\subsection{Phase 2 Attention Analysis}
% ============================================================

Attention was analyzed using heatmaps, entropy, and mean attention
distance.

\begin{table}[H]
\centering
\small
\caption{Phase 2 attention statistics.}
\begin{tabular}{lcc}
\toprule
 & \textbf{Telugu} & \textbf{Nepali}\\
\midrule
Mean attention entropy & 2.187 & 2.087\\
Mean attention distance & 11.78 & 12.73\\
Early-layer distance & 15.37 & 14.76\\
Late-layer distance & 16.25 & 17.42\\
\bottomrule
\end{tabular}
\end{table}

The following figures show multiple heads instead of only one
representative head. Four heads from the early layer and four heads
from the late layer are shown for each model.

\begin{figure}[H]
\centering

\begin{subfigure}{0.235\linewidth}
    \centering
    \safeimage{figures/telugu_attention_early_layer_00_head_00.png}
    \caption{L0 H0}
\end{subfigure}
\hfill
\begin{subfigure}{0.235\linewidth}
    \centering
    \safeimage{figures/telugu_attention_early_layer_00_head_01.png}
    \caption{L0 H1}
\end{subfigure}
\hfill
\begin{subfigure}{0.235\linewidth}
    \centering
    \safeimage{figures/telugu_attention_early_layer_00_head_02.png}
    \caption{L0 H2}
\end{subfigure}
\hfill
\begin{subfigure}{0.235\linewidth}
    \centering
    \safeimage{figures/telugu_attention_early_layer_00_head_03.png}
    \caption{L0 H3}
\end{subfigure}

\vspace{2pt}

\begin{subfigure}{0.235\linewidth}
    \centering
    \safeimage{figures/telugu_attention_late_layer_09_head_00.png}
    \caption{L9 H0}
\end{subfigure}
\hfill
\begin{subfigure}{0.235\linewidth}
    \centering
    \safeimage{figures/telugu_attention_late_layer_09_head_01.png}
    \caption{L9 H1}
\end{subfigure}
\hfill
\begin{subfigure}{0.235\linewidth}
    \centering
    \safeimage{figures/telugu_attention_late_layer_09_head_02.png}
    \caption{L9 H2}
\end{subfigure}
\hfill
\begin{subfigure}{0.235\linewidth}
    \centering
    \safeimage{figures/telugu_attention_late_layer_09_head_03.png}
    \caption{L9 H3}
\end{subfigure}

\caption{Telugu Phase 2 attention maps. The top row shows four heads
from the early layer and the bottom row shows four heads from the
late layer. Axes are query position and key position, with color
representing attention weight.}
\end{figure}

\begin{figure}[H]
\centering

\begin{subfigure}{0.235\linewidth}
    \centering
    \safeimage{figures/nepali_attention_early_layer_00_head_00.png}
    \caption{L0 H0}
\end{subfigure}
\hfill
\begin{subfigure}{0.235\linewidth}
    \centering
    \safeimage{figures/nepali_attention_early_layer_00_head_01.png}
    \caption{L0 H1}
\end{subfigure}
\hfill
\begin{subfigure}{0.235\linewidth}
    \centering
    \safeimage{figures/nepali_attention_early_layer_00_head_02.png}
    \caption{L0 H2}
\end{subfigure}
\hfill
\begin{subfigure}{0.235\linewidth}
    \centering
    \safeimage{figures/nepali_attention_early_layer_00_head_03.png}
    \caption{L0 H3}
\end{subfigure}

\vspace{2pt}

\begin{subfigure}{0.235\linewidth}
    \centering
    \safeimage{figures/nepali_attention_late_layer_10_head_00.png}
    \caption{L10 H0}
\end{subfigure}
\hfill
\begin{subfigure}{0.235\linewidth}
    \centering
    \safeimage{figures/nepali_attention_late_layer_10_head_01.png}
    \caption{L10 H1}
\end{subfigure}
\hfill
\begin{subfigure}{0.235\linewidth}
    \centering
    \safeimage{figures/nepali_attention_late_layer_10_head_02.png}
    \caption{L10 H2}
\end{subfigure}
\hfill
\begin{subfigure}{0.235\linewidth}
    \centering
    \safeimage{figures/nepali_attention_late_layer_10_head_03.png}
    \caption{L10 H3}
\end{subfigure}

\caption{Nepali Phase 2 attention maps. The top row shows four heads
from the early layer and the bottom row shows four heads from the
late layer.}
\end{figure}

The attention maps show that different heads behave differently. Some
heads focus more on nearby positions, while others attend to positions
farther away. Some heads also appear to attend based on the content
of the tokens rather than only their position.

One interesting pattern is that the attention distance does not
simply increase as the model gets deeper. Both models show broader
attention in the early layers, more local attention in the middle
layers, and broader attention again in the later layers. For Telugu,
the mean distance goes from 15.37 in the first layer down to 4.34
around the middle layers, and then goes back up to 16.25 in the final
layer. Nepali shows a similar pattern, going from 14.76 to 4.52 and
then increasing to 17.42 in the final layer.

There are also clear examples of individual head specialization.
For Telugu, Layer 5 Head 1 has very low entropy (0.097) and a mean
distance of about 1.04, making it a strong local example. For Nepali,
Layer 5 Head 5 has entropy 0.081 and mean distance 1.01. On the
other side, Nepali Layer 7 Head 6 has a mean distance of 24.54, while
Layer 7 Head 2 has a distance of 23.26. These show that the heads are
not all using the context in the same way.

Overall, the two models show similar types of head specialization,
although Nepali has a slightly higher overall mean attention distance
(12.73 vs. 11.78) and a higher late-layer distance (17.42 vs. 16.25).

% ============================================================
\section{Phase 3: Reasoning Fine-tuning}
% ============================================================

\subsection{Synthetic Reasoning Dataset}

Each language uses a generated synthetic reasoning dataset.
The main dataset contains 20,000 training, 2,000 validation, and
3,000 test examples.

The tasks include direct comparison, reverse comparison, equality,
two-hop transitivity, three-hop transitivity, and ordering. Attributes
include height, age, weight, price, and quantity.

Examples are generated independently for each language and use the
target language's script and phrasing. Entity names, numerical
values, relation directions, and statement ordering are varied to
reduce template memorization.

The generator has explicit leakage checks. Exact prompts must be
unique across train, validation, and test. Test templates are kept
separate from train/validation templates. Entity names are allowed to
occur in different splits, since the same reasoning task should still
work with different entities. The original train, validation, and
test files are not modified during the experiments.

\subsection{Fine-tuning Results}

The primary Phase 3 comparison uses the pure synthetic reasoning
condition containing 20K reasoning examples and 0K original
pretraining examples.

\begin{table}[H]
\centering
\small
\caption{Reasoning performance before and after fine-tuning.}
\begin{tabular}{lcc}
\toprule
 & \textbf{Telugu} & \textbf{Nepali}\\
\midrule
Test examples & 3,000 & 3,000\\
Pretrained exact match & 0.00\% & 0.00\%\\
Fine-tuned exact match & 42.10\% & 13.90\%\\
Absolute improvement & 42.10 pp & 13.90 pp\\
Pretrained answer-token accuracy & 6.75\% & 1.29\%\\
Fine-tuned answer-token accuracy & 45.60\% & 43.88\%\\
\bottomrule
\end{tabular}
\end{table}

Fine-tuning showed a good improvement in both models. Telugu
reaches 42.10\% exact match, while Nepali reaches 13.90\%.

The answer-token results show additional information about the Nepali
model. Its answer-token accuracy rises from 1.29\% to 43.88\%, even
though exact-match accuracy reaches only 13.90\%. Therefore, many
Nepali outputs contain the expected answer token but do not match
exactly.

For Nepali, reasoning difficulty also varies by task. Equality
examples show almost 100\% accuracy, while direct comparison, reverse
comparison, transitive-2, transitive-3, and ordering remain much more
difficult. Difficulty also has a large effect: easy examples reach
27.00\% exact match compared with 3.68\% for hard examples.

\subsection{Qualitative Reasoning Examples}

The following examples are taken from the actual Phase 3 evaluation
outputs.

\begin{table}[H]
\centering
\scriptsize
\begin{tabular}{p{0.25\linewidth}p{0.69\linewidth}}

\toprule

\multicolumn{2}{c}{\textbf{Telugu}}\\

\midrule

\textbf{Successful example} &
\textbf{Prompt:} {\telugufont
దినేష్ యొక్క వయస్సు శశాంక్ కంటే ఎక్కువగా ఉంది.

సందీప్ యొక్క వయస్సు దినేష్ కంటే ఎక్కువగా ఉంది. దినేష్, సందీప్,
శశాంక్ -- వీరిలో అత్యధిక వయస్సు ఉన్నది ఎవరు?
}

\textbf{Gold:} {\telugufont సందీప్}

\textbf{Fine-tuned:} {\telugufont సందీప్}

\\[10pt]

\textbf{Failure example} &
\textbf{Prompt:} {\telugufont
విశాల్ యొక్క వయస్సుతో పోలిస్తే అశోక్ యొక్క
వయస్సు ఎక్కువగా ఉంది. శ్రీనివాస్ యొక్క వయస్సు విశాల్ కంటే తక్కువ.

శ్రీనివాస్, విశాల్, అశోక్ -- వీరిలో అత్యధిక వయస్సు ఉన్నది ఎవరు?
}

\textbf{Gold:} {\telugufont అశోక్}

\textbf{Fine-tuned:} {\telugufont విశాల్}

\\

\midrule

\multicolumn{2}{c}{\textbf{Nepali}}\\

\midrule

\textbf{Successful example} &
\textbf{Prompt:} {\devanagarifont
विक्रम र दिपेशको तौल उस्तै छ।

विक्रम र दिपेशको तौल बराबर छ?
}

\textbf{Gold:} {\devanagarifont हो}

\textbf{Fine-tuned:} {\devanagarifont हो}

\\[10pt]

\textbf{Failure example} &
\textbf{Prompt:} {\devanagarifont
मनोजभन्दा नविनको उचाइ धेरै देखिन्छ। मनोजभन्दा
बिमलको उचाइ थोरै देखिन्छ। बिमलको उचाइ बिराजको तुलनामा बढी रहेको छ।

बिराज, बिमल, नविन, मनोज मा न्यूनतम उचाइ कसको छ?
}

\textbf{Gold:} {\devanagarifont बिराज}

\textbf{Fine-tuned:} {\devanagarifont बिमल}

\\

\bottomrule

\end{tabular}

\caption{Examples of successful and failed reasoning predictions
after fine-tuning.}
\end{table}

The successful examples show that the fine-tuned models can produce
the required short answer for some comparison tasks. The failures
show that the models can still make mistakes when several relations
have to be combined, especially in transitive or harder ordering
examples.

\subsection{Reasoning/Data Mixture Experiment}

Additional experiments varied the amount of synthetic reasoning data
while keeping the total fine-tuning set at 20K examples.

\begin{table}[H]
\centering
\scriptsize
\caption{Reasoning-data mixture results. Values are exact-match
accuracy on the same 3,000-example test set.}
\begin{tabular}{lcc}
\toprule
Reasoning / Original data & Telugu & Nepali\\
\midrule
2K / 18K & 31.07\% & 5.67\%\\
5K / 15K & 36.37\% & 10.67\%\\
8K / 12K & 37.30\% & 11.07\%\\
12K / 8K & 36.67\% & 14.60\%\\
16K / 4K & 36.47\% & 15.10\%\\
20K / 0K & 42.10\% & 13.90\%\\
\bottomrule
\end{tabular}
\end{table}

The relationship between the reasoning-data fraction and accuracy is
not strictly monotonic for every intermediate mixture.

Telugu obtains its highest measured exact-match accuracy in the pure
reasoning condition. Nepali reaches its highest measured mixture
accuracy of 15.10\% with 16K reasoning and 4K original-language
examples, compared with 13.90\% in the pure reasoning condition.

This experiment shows that the fine-tuning objective matters in
addition to the number of reasoning examples. Adding original
language-modeling examples changes the distribution of the fine-tuning
objective and therefore changes the resulting task performance.

% ============================================================
\subsection{Post-Fine-tuning Attention Analysis}
% ============================================================

Attention was re-evaluated on reasoning prompts using the pretrained
and fine-tuned checkpoints.

\begin{table}[H]
\centering
\small
\caption{Attention statistics before and after reasoning fine-tuning.}
\begin{tabular}{llrrrr}
\toprule
Model & State &
Entropy & Distance &
$\Delta$ Entropy & $\Delta$ Distance\\
\midrule
Telugu & Pretrained
& 2.040 & 9.442 & -- & --\\
Telugu & Fine-tuned
& 1.894 & 8.891 & -0.147 & -0.551\\
Nepali & Pretrained
& 1.843 & 9.579 & -- & --\\
Nepali & Fine-tuned
& 1.827 & 8.632 & -0.016 & -0.947\\
\bottomrule
\end{tabular}
\end{table}

The pretrained and fine-tuned attention maps were compared directly
using the same reasoning prompts.

\begin{figure}[H]
\centering

\begin{subfigure}{0.235\linewidth}
    \centering
    \safeimage{figures/telugu_attention_pretrained_early_layer.png}
    \caption{Telugu pretrained early}
\end{subfigure}
\hfill
\begin{subfigure}{0.235\linewidth}
    \centering
    \safeimage{figures/telugu_attention_finetuned_early_layer.png}
    \caption{Telugu fine-tuned early}
\end{subfigure}
\hfill
\begin{subfigure}{0.235\linewidth}
    \centering
    \safeimage{figures/telugu_attention_pretrained_late_layer.png}
    \caption{Telugu pretrained late}
\end{subfigure}
\hfill
\begin{subfigure}{0.235\linewidth}
    \centering
    \safeimage{figures/telugu_attention_finetuned_late_layer.png}
    \caption{Telugu fine-tuned late}
\end{subfigure}

\caption{Telugu attention before and after reasoning fine-tuning.
The same early/late layers and representative reasoning setup are
used for the comparison.}
\end{figure}

\begin{figure}[H]
\centering

\begin{subfigure}{0.235\linewidth}
    \centering
    \safeimage{figures/nepali_attention_pretrained_early_layer.png}
    \caption{Nepali pretrained early}
\end{subfigure}
\hfill
\begin{subfigure}{0.235\linewidth}
    \centering
    \safeimage{figures/nepali_attention_finetuned_early_layer.png}
    \caption{Nepali fine-tuned early}
\end{subfigure}
\hfill
\begin{subfigure}{0.235\linewidth}
    \centering
    \safeimage{figures/nepali_attention_pretrained_late_layer.png}
    \caption{Nepali pretrained late}
\end{subfigure}
\hfill
\begin{subfigure}{0.235\linewidth}
    \centering
    \safeimage{figures/nepali_attention_finetuned_late_layer.png}
    \caption{Nepali fine-tuned late}
\end{subfigure}

\caption{Nepali attention before and after reasoning fine-tuning.}
\end{figure}

Fine-tuning decreases average attention entropy for both models,
which means the attention becomes more concentrated on the reasoning
prompts. Mean attention distance also decreases, by approximately
0.55 tokens for Telugu and 0.95 tokens for Nepali.

The change is larger for Nepali in terms of mean distance, while the
entropy change is much smaller. These are aggregate changes, so they
do not mean that every individual head became more local. Individual
heads can still behave differently after fine-tuning.

% ============================================================
\section{Final Comparison and Discussion}
% ============================================================

\begin{table}[H]
\centering
\small
\caption{Consolidated results for the two resource tiers.}
\begin{tabular}{p{0.36\linewidth}cc}
\toprule
 & \textbf{Model H: Telugu} & \textbf{Model L: Nepali}\\
\midrule
Parameters & 23.04M & 25.95M\\
Training tokens & 486.03M & 500M\\
Vocabulary & 6K & 8K\\
Manual fraction & 25.00\% & 18.24\%\\
Tokenizer fertility & 2.242 & 1.667\\
Characters/token & 3.53 & 3.88\\
Test PPL & 22.73 & 30.23\\
Test BPB & 0.467 & 0.846\\
Greedy BLEU-4 & 8.14 & 3.43\\
Greedy repetition & 54.3\% & 71.2\%\\
Greedy Distinct-2 & .405 & .226\\
Reasoning exact match & 42.10\% & 13.90\%\\
\bottomrule
\end{tabular}
\end{table}

\subsection{1. How did data scale and quality differ?}

The final training corpora are similar in token count, with 486.03M
tokens for Telugu and 500M for Nepali. So the raw number of training
tokens alone does not explain the difference in the results.

The main difference is in the availability and composition of manual
data. Telugu had enough manually collected material to use a 25\%
manual fraction, while Nepali could use only 18.24\% because the
available manual corpus was smaller.

There were also differences within the Nepali corpus itself. Its OCR
material has much higher tokenizer fertility than its main web
sources. This suggests that the quality of the training data was not
the same across all sources.

\subsection{2. How did language-modeling and reasoning results compare?}

Telugu obtains lower test cross-entropy, perplexity, and BPB. Its
greedy generations also have lower repetition and higher diversity
than Nepali.

After synthetic reasoning fine-tuning, both models improve from 0\%
pretrained exact match. Telugu reaches 42.10\%, while Nepali reaches
13.90\%. The answer-token accuracy gives some additional information
for Nepali. It learns a lot of the required answer behavior, but many
outputs still do not match the expected answer exactly.

So, the difference between the two models can be seen in both the
language-modeling results and the synthetic reasoning results.

\subsection{3. Which tokenizer and corpus factors affected Model L?}

The main corpus-related limitation for Nepali was the smaller amount
of manually collected data, which prevented the desired 20\% manual
fraction from being reached.

The tokenizer statistics also show that the Nepali corpus is quite
heterogeneous. The OCR material produces much more fragmented
tokenization than the main web sources. This means that some parts of
the training corpus are harder to tokenize cleanly than others.

The higher test BPB of 0.846 compared with Telugu's 0.467 also shows
that the difference is present beyond just the choice of subword
vocabulary.

\subsection{4. What evidence explains the observed differences?}

Several results point to the same general differences:

\begin{itemize}
    \item Telugu has lower held-out cross-entropy and perplexity.
    \item Telugu has substantially lower BPB.
    \item Nepali greedy generations have a higher repetition rate.
    \item Nepali has a smaller achievable manual-data fraction.
    \item Nepali contains source subsets with substantially higher
          tokenizer fertility, especially OCR material.
    \item Telugu obtains higher exact-match reasoning accuracy after
          fine-tuning.
    \item Both models show specialized local and long-range attention
          heads.
    \item Fine-tuning changes attention entropy and distance on
          reasoning prompts.
    \item Removing positional information slightly worsens Telugu
          intrinsic language-modeling performance.
\end{itemize}

These results suggest that corpus availability, source quality,
tokenizer behavior, and resource level all have a role in the
differences between the models. However, this is not a controlled
causal experiment. The two languages also differ in script,
linguistic structure, corpus composition, and tokenizer configuration.
So the observed difference should not be attributed to just one
factor.

% ============================================================

\section{Limitations}

% ============================================================

The two models cannot be treated as identical experiments differing
only in resource level. Their corpora, languages, scripts, vocabulary
sizes, and manually collected data differ.

Perplexity is tokenizer-dependent, making BPB more useful for
cross-language intrinsic comparison.

The reasoning evaluation uses a controlled synthetic benchmark, so its
accuracy should not be interpreted as a general measure of reasoning
ability. Exact match is also stricter than answer-token accuracy since
additional incorrect output makes an otherwise correct answer count as
incorrect.

The attention analysis uses representative held-out sequences and
should therefore be interpreted as evidence of attention patterns,
rather than proof that a particular head always performs a specific
linguistic function.

% ============================================================

\section{Conclusion}

% ============================================================

The project produced two independent approximately 25M-parameter
decoder-only Transformer language models trained from scratch for
Telugu and Nepali, with independent datasets, tokenizers, vocabularies,
and weights.

The final Telugu model achieves test perplexity 22.73 and BPB 0.467,
while the Nepali model achieves perplexity 30.23 and BPB 0.846.
Generation experiments show the expected trade-off between sampling
temperature, reference overlap, repetition, and diversity.

Synthetic reasoning fine-tuning improves both models. Exact-match
accuracy increases from 0\% to 42.10\% for Telugu and from 0\% to
13.90\% for Nepali in the pure reasoning condition. Attention analysis
shows different interaction ranges across heads, while post-fine-tuning
statistics show somewhat more concentrated attention on reasoning
prompts.

The positional-embedding ablation shows a small degradation in
intrinsic language-modeling performance when positional information is
removed, although some generation metrics move in the other direction.

Overall, similar pretraining token counts do not necessarily give the
same training conditions. Manual data, source quality, tokenizer
behavior, and language characteristics are relevant to the differences
between the two models.

\end{document}